\documentclass[10pt,a4paper]{amsart}
\usepackage{amsmath,amssymb,amsthm,graphicx}
\usepackage{verbatim}
\newenvironment{reference}{\comment}{\endcomment}
%\newenvironment{reference}{}{}
\newenvironment{slogan}{\comment}{\endcomment}
\newenvironment{history}{\comment}{\endcomment}
% For commutative diagrams we use Xy-pic
\usepackage[all]{xy}

% We use 2cell for 2-commutative diagrams.
\xyoption{2cell}
% Theorem environments.
%
\theoremstyle{plain}
\newtheorem{theorem}[subsection]{Theorem}
\newtheorem{proposition}[subsection]{Proposition}
\newtheorem{lemma}[subsection]{Lemma}

\theoremstyle{definition}
\newtheorem{definition}[subsection]{Definition}
\newtheorem{example}[subsection]{Example}
\newtheorem{exercise}[subsection]{Exercise}
\newtheorem{situation}[subsection]{Situation}

\theoremstyle{remark}
\newtheorem{remark}[subsection]{Remark}
\newtheorem{remarks}[subsection]{Remarks}

\numberwithin{equation}{subsection}

% Macros
%
\def\lim{\mathop{\mathrm{lim}}\nolimits}
\def\colim{\mathop{\mathrm{colim}}\nolimits}
\def\Spec{\mathop{\mathrm{Spec}}}
\def\Hom{\mathop{\mathrm{Hom}}\nolimits}
\def\Ext{\mathop{\mathrm{Ext}}\nolimits}
\def\SheafHom{\mathop{\mathcal{H}\!\mathit{om}}\nolimits}
\def\SheafExt{\mathop{\mathcal{E}\!\mathit{xt}}\nolimits}
\def\Sch{\mathit{Sch}}
\def\Mor{\mathop{\mathrm{Mor}}\nolimits}
\def\Ob{\mathop{\mathrm{Ob}}\nolimits}
\def\Sh{\mathop{\mathit{Sh}}\nolimits}
\def\NL{\mathop{N\!L}\nolimits}
\def\CH{\mathop{\mathrm{CH}}\nolimits}
\def\proetale{{pro\text{-}\acute{e}tale}}
\def\etale{{\acute{e}tale}}
\def\QCoh{\mathit{QCoh}}
\def\Ker{\mathop{\mathrm{Ker}}}
\def\Im{\mathop{\mathrm{Im}}}
\def\Coker{\mathop{\mathrm{Coker}}}
\def\Coim{\mathop{\mathrm{Coim}}}

% Boxtimes
%
\DeclareMathSymbol{\boxtimes}{\mathbin}{AMSa}{"02}

%
% Macros for moduli stacks/spaces
%
\def\QCohstack{\mathcal{QC}\!\mathit{oh}}
\def\Cohstack{\mathcal{C}\!\mathit{oh}}
\def\Spacesstack{\mathcal{S}\!\mathit{paces}}
\def\Quotfunctor{\mathrm{Quot}}
\def\Hilbfunctor{\mathrm{Hilb}}
\def\Curvesstack{\mathcal{C}\!\mathit{urves}}
\def\Polarizedstack{\mathcal{P}\!\mathit{olarized}}
\def\Complexesstack{\mathcal{C}\!\mathit{omplexes}}
% \Pic is the operator that assigns to X its picard group, usage \Pic(X)
% \Picardstack_{X/B} denotes the Picard stack of X over B
% \Picardfunctor_{X/B} denotes the Picard functor of X over B
\def\Pic{\mathop{\mathrm{Pic}}\nolimits}
\def\Picardstack{\mathcal{P}\!\mathit{ic}}
\def\Picardfunctor{\mathrm{Pic}}
\def\Deformationcategory{\mathcal{D}\!\mathit{ef}}
\usepackage{cleveref}
\crefname{originallemma}{original lemma}{original lemmas}
\Crefname{originallemma}{Original lemma}{Original lemmas}
\makeatletter
\newlabel{coherent-lemma-surjective-finite-morphism-ample}{{0B5V}{source}{}{sourcepointer-1}{}}
\newlabel{coherent-lemma-surjective-finite-morphism-ample@cref}{{[originallemma][0][]0B5V}{[1][0][]source}{}{}{}}
\makeatother
\crefname{originallemma}{original lemma}{original lemmas}
\Crefname{originallemma}{Original lemma}{Original lemmas}
\makeatletter
\newlabel{lemma-affine-morphism-projection-ideal}{{0GF8}{source}{}{sourcepointer-2}{}}
\newlabel{lemma-affine-morphism-projection-ideal@cref}{{[originallemma][0][]0GF8}{[1][0][]source}{}{}{}}
\makeatother
\crefname{originallemma}{original lemma}{original lemmas}
\Crefname{originallemma}{Original lemma}{Original lemmas}
\makeatletter
\newlabel{lemma-finite-higher-direct-image-zero}{{0A4K}{source}{}{sourcepointer-3}{}}
\newlabel{lemma-finite-higher-direct-image-zero@cref}{{[originallemma][0][]0A4K}{[1][0][]source}{}{}{}}
\makeatother
\crefname{originallemma}{original lemma}{original lemmas}
\Crefname{originallemma}{Original lemma}{Original lemmas}
\makeatletter
\newlabel{lemma-finite-morphism-Noetherian}{{0GF7}{source}{}{sourcepointer-4}{}}
\newlabel{lemma-finite-morphism-Noetherian@cref}{{[originallemma][0][]0GF7}{[1][0][]source}{}{}{}}
\makeatother
\crefname{originallemma}{original lemma}{original lemmas}
\Crefname{originallemma}{Original lemma}{Original lemmas}
\makeatletter
\newlabel{lemma-finite-pushforward-coherent}{{07UH}{source}{}{sourcepointer-5}{}}
\newlabel{lemma-finite-pushforward-coherent@cref}{{[originallemma][0][]07UH}{[1][0][]source}{}{}{}}
\makeatother
\crefname{originallemma}{original lemma}{original lemmas}
\Crefname{originallemma}{Original lemma}{Original lemmas}
\makeatletter
\newlabel{lemma-i-star-equivalence}{{08AM}{source}{}{sourcepointer-6}{}}
\newlabel{lemma-i-star-equivalence@cref}{{[originallemma][0][]08AM}{[1][0][]source}{}{}{}}
\makeatother
\crefname{originallemma}{original lemma}{original lemmas}
\Crefname{originallemma}{Original lemma}{Original lemmas}
\makeatletter
\newlabel{lemma-property-higher-rank-cohomological}{{07UT}{source}{}{sourcepointer-7}{}}
\newlabel{lemma-property-higher-rank-cohomological@cref}{{[originallemma][0][]07UT}{[1][0][]source}{}{}{}}
\makeatother
\crefname{originallemma}{original lemma}{original lemmas}
\Crefname{originallemma}{Original lemma}{Original lemmas}
\makeatletter
\newlabel{lemma-vanshing-gives-ample}{{0GFA}{source}{}{sourcepointer-8}{}}
\newlabel{lemma-vanshing-gives-ample@cref}{{[originallemma][0][]0GFA}{[1][0][]source}{}{}{}}
\makeatother
\crefname{originalsection}{original section}{original sections}
\Crefname{originalsection}{Original section}{Original sections}
\makeatletter
\newlabel{sites-cohomology-section-leray}{{072X}{source}{}{sourcepointer-9}{}}
\newlabel{sites-cohomology-section-leray@cref}{{[originalsection][0][]072X}{[1][0][]source}{}{}{}}
\makeatother
\crefname{originalsection}{original section}{original sections}
\Crefname{originalsection}{Original section}{Original sections}
\makeatletter
\newlabel{sites-cohomology-section-projection-formula}{{0943}{source}{}{sourcepointer-10}{}}
\newlabel{sites-cohomology-section-projection-formula@cref}{{[originalsection][0][]0943}{[1][0][]source}{}{}{}}
\makeatother
\crefname{originallemma}{original lemma}{original lemmas}
\Crefname{originallemma}{Original lemma}{Original lemmas}
\makeatletter
\newlabel{spaces-morphisms-lemma-integral-local}{{03ZQ}{source}{}{sourcepointer-11}{}}
\newlabel{spaces-morphisms-lemma-integral-local@cref}{{[originallemma][0][]03ZQ}{[1][0][]source}{}{}{}}
\makeatother
\title{Ampleness of an invertible sheaf passes up and down along a finite surjective morphism of algebraic spaces proper over a Noetherian ring}
\author{The Stacks Project Authors}
\date{stacks-project:spaces-cohomology.tex:local-snapshot-20261012; GNU FDL 1.2 or later}
\begin{document}
\maketitle
\noindent\textit{Editorial note.} All delivered mathematical text is the author's own, copied from the pinned Stacks Project chapter file spaces-cohomology.tex; no statement or proof step was rewritten, reordered or removed.\par\medskip
\noindent\textit{Editorial note.} Preamble provenance: the pinned chapter file opens with \textbackslash{}input\{preamble\} and carries no preamble text of its own, so the project's single Stacks-wide preamble staged at sources/stacks-spaces-cohomology/preamble.tex is excerpted verbatim as three spans, lines 8-12 (verbatim package and the project's own reference/slogan/history display environments), lines 14-18 (Xy-pic setup) and lines 154-217 (theorem declarations and mathematical macros). The wrapper class and the amsmath/amssymb/amsthm/graphicx base packages are supplied by the assembler. No macro was invented or redefined.\par\medskip
\noindent\textit{Editorial note.} Same-chapter pointers are located at their \textbackslash{}label lines inside the pinned chapter file; cross-chapter pointers are located in the project's own tag list staged in the snapshot (sources/stacks-chow/tags.txt), at the verbatim lines carrying their canonical tags.\par\medskip
\noindent\textit{Allowed standard prerequisites (editorial).}\begin{itemize}
\item A finite morphism of algebraic spaces is representable by schemes. Morphisms of Spaces, label spaces-morphisms-lemma-integral-local (invoked at proof line 4147). It shows that if \$X\$ is a scheme then so is \$Y\$, which is the scheme-theoretic half of the implication (1) \$\textbackslash{}Rightarrow\$ (2).
\item The scheme case of the equivalence for a finite surjective morphism. Cohomology of Schemes, label coherent-lemma-surjective-finite-morphism-ample (invoked at proof line 4149). It supplies ampleness of \$f\textasciicircum{}*\textbackslash{}mathcal\{L\}\$ once \$X\$ is a scheme and \$\textbackslash{}mathcal\{L\}\$ is ample.
\item The cohomological criterion for ampleness, in both directions. Same chapter spaces-cohomology.tex, label lemma-vanshing-gives-ample (invoked at proof lines 4158 and 4191). It converts the property \$P\$ (vanishing of the cohomology of all twists of every coherent module) into ampleness of \$\textbackslash{}mathcal\{L\}\$, and over the scheme \$Y\$ it produces the integer \$n\_0\$ for the sheaf \$i'\_*\textbackslash{}mathcal\{I\}'\$ twisted by \$f\textasciicircum{}*\textbackslash{}mathcal\{L\}\$.
\item The criterion reducing a property of coherent sheaves on a Noetherian algebraic space to irreducible supports. Same chapter spaces-cohomology.tex, label lemma-property-higher-rank-cohomological (invoked at proof lines 4159 and 4171). Conditions (1) and (2) are verified directly for \$P\$, and condition (3) is met by the sheaf \$\textbackslash{}mathcal\{G\} = f'\_*\textbackslash{}mathcal\{O\}\_\{Z'\}\$ on every reduced irreducible closed subspace.
\item The auxiliary finite morphism over an irreducible closed subspace, and the projection formula for an affine morphism. Same chapter spaces-cohomology.tex, labels lemma-finite-morphism-Noetherian (invoked at proof lines 4169 and 4173) and lemma-affine-morphism-projection-ideal (invoked at proof line 4179). The first provides \$Z'\$ with the finite morphisms \$i' : Z' \textbackslash{}to Y\$ and \$f' : Z' \textbackslash{}to Z\$ and the support property (3)(a); the second gives \$f'\_*\textbackslash{}mathcal\{I\}' = \textbackslash{}mathcal\{I\}\textbackslash{}mathcal\{G\}\$, which identifies the subsheaf \$\textbackslash{}mathcal\{G\}'\$ used in (3)(b).
\item Coherence of the relevant pushforwards and vanishing of the higher direct images of the finite morphism. Same chapter spaces-cohomology.tex, labels lemma-finite-pushforward-coherent (invoked at proof line 4188), lemma-i-star-equivalence (invoked at proof line 4189) and lemma-finite-higher-direct-image-zero (invoked at proof line 4220), together with the projection formula and the Leray spectral sequence of Cohomology on Sites (proof lines 4216-4219). They let the cohomology of \$i\_*\textbackslash{}mathcal\{G\}' \textbackslash{}otimes \textbackslash{}mathcal\{L\}\textasciicircum{}\{\textbackslash{}otimes n\}\$ on \$X\$ be rewritten as the cohomology of \$i'\_*\textbackslash{}mathcal\{I\}' \textbackslash{}otimes f\textasciicircum{}*\textbackslash{}mathcal\{L\}\textasciicircum{}\{\textbackslash{}otimes n\}\$ on the scheme \$Y\$, where the vanishing for \$n \textbackslash{}geq n\_0\$ is known.
\end{itemize}\medskip
\ifdefined\hypertarget\hypertarget{sourcepointer-1}{}\fi
\noindent\textit{Source reference.} Original lemma 0B5V: sources/stacks-chow/tags.txt, source lines 14048--14048. Navigation pointer for \textbackslash{}ref\{coherent-lemma-surjective-finite-morphism-ample\} used inside the counted span (line 4150); the argument that consumes it is retained in the span.\par\medskip
\ifdefined\hypertarget\hypertarget{sourcepointer-2}{}\fi
\noindent\textit{Source reference.} Original lemma 0GF8: sources/stacks-spaces-cohomology/pinned-original-source.tex, source lines 3386--3386. Navigation pointer for \textbackslash{}ref\{lemma-affine-morphism-projection-ideal\} used inside the counted span (line 4179); the argument that consumes it is retained in the span.\par\medskip
\ifdefined\hypertarget\hypertarget{sourcepointer-3}{}\fi
\noindent\textit{Source reference.} Original lemma 0A4K: sources/stacks-spaces-cohomology/pinned-original-source.tex, source lines 496--496. Navigation pointer for \textbackslash{}ref\{lemma-finite-higher-direct-image-zero\} used inside the counted span (line 4220); the argument that consumes it is retained in the span.\par\medskip
\ifdefined\hypertarget\hypertarget{sourcepointer-4}{}\fi
\noindent\textit{Source reference.} Original lemma 0GF7: sources/stacks-spaces-cohomology/pinned-original-source.tex, source lines 3351--3351. Navigation pointer for \textbackslash{}ref\{lemma-finite-morphism-Noetherian\} used inside the counted span (line 4169); the argument that consumes it is retained in the span.\par\medskip
\ifdefined\hypertarget\hypertarget{sourcepointer-5}{}\fi
\noindent\textit{Source reference.} Original lemma 07UH: sources/stacks-spaces-cohomology/pinned-original-source.tex, source lines 2284--2284. Navigation pointer for \textbackslash{}ref\{lemma-finite-pushforward-coherent\} used inside the counted span (line 4188); the argument that consumes it is retained in the span.\par\medskip
\ifdefined\hypertarget\hypertarget{sourcepointer-6}{}\fi
\noindent\textit{Source reference.} Original lemma 08AM: sources/stacks-spaces-cohomology/pinned-original-source.tex, source lines 2259--2259. Navigation pointer for \textbackslash{}ref\{lemma-i-star-equivalence\} used inside the counted span (line 4189); the argument that consumes it is retained in the span.\par\medskip
\ifdefined\hypertarget\hypertarget{sourcepointer-7}{}\fi
\noindent\textit{Source reference.} Original lemma 07UT: sources/stacks-spaces-cohomology/pinned-original-source.tex, source lines 2652--2652. Navigation pointer for \textbackslash{}ref\{lemma-property-higher-rank-cohomological\} used inside the counted span (line 4159); the argument that consumes it is retained in the span.\par\medskip
\ifdefined\hypertarget\hypertarget{sourcepointer-8}{}\fi
\noindent\textit{Source reference.} Original lemma 0GFA: sources/stacks-spaces-cohomology/pinned-original-source.tex, source lines 4109--4109. Navigation pointer for \textbackslash{}ref\{lemma-vanshing-gives-ample\} used inside the counted span (line 4158); the argument that consumes it is retained in the span.\par\medskip
\ifdefined\hypertarget\hypertarget{sourcepointer-9}{}\fi
\noindent\textit{Source reference.} Original section 072X: sources/stacks-chow/tags.txt, source lines 9041--9041. Navigation pointer for \textbackslash{}ref\{sites-cohomology-section-leray\} used inside the counted span (line 4219); the argument that consumes it is retained in the span.\par\medskip
\ifdefined\hypertarget\hypertarget{sourcepointer-10}{}\fi
\noindent\textit{Source reference.} Original section 0943: sources/stacks-chow/tags.txt, source lines 11536--11536. Navigation pointer for \textbackslash{}ref\{sites-cohomology-section-projection-formula\} used inside the counted span (line 4218); the argument that consumes it is retained in the span.\par\medskip
\ifdefined\hypertarget\hypertarget{sourcepointer-11}{}\fi
\noindent\textit{Source reference.} Original lemma 03ZQ: sources/stacks-chow/tags.txt, source lines 5229--5229. Navigation pointer for \textbackslash{}ref\{spaces-morphisms-lemma-integral-local\} used inside the counted span (line 4148); the argument that consumes it is retained in the span.\par\medskip
\begin{lemma}
\label{lemma-surjective-finite-morphism-ample}
Let $R$ be a Noetherian ring. Let $f : Y \to X$ be a morphism of
algebraic spaces proper over $R$. Let $\mathcal{L}$ be an
invertible $\mathcal{O}_X$-module. Assume $f$ is finite and surjective.
The following are equivalent
\begin{enumerate}
\item $X$ is a scheme and $\mathcal{L}$ is ample, and
\item $Y$ is a scheme and $f^*\mathcal{L}$ is ample.
\end{enumerate}
\end{lemma}

\begin{proof}
Assume (1). Then $Y$ is a scheme as a finite morphism is representable
(by schemes), see Morphisms of Spaces, Lemma
\ref{spaces-morphisms-lemma-integral-local}. Hence (2)
follows from Cohomology of Schemes, Lemma
\ref{coherent-lemma-surjective-finite-morphism-ample}.

\medskip\noindent
Assume (2). Let $P$ be the following property on
coherent $\mathcal{O}_X$-modules $\mathcal{F}$: there exists an $n_0$
such that $H^p(X, \mathcal{F} \otimes \mathcal{L}^{\otimes n}) = 0$
for all $n \geq n_0$ and $p > 0$. We will prove that $P$ holds
for any coherent $\mathcal{O}_X$-module $\mathcal{F}$, which implies
$\mathcal{L}$ is ample by Lemma \ref{lemma-vanshing-gives-ample}.
We are going to apply Lemma \ref{lemma-property-higher-rank-cohomological}.
Thus we have to verify (1), (2) and (3) of that lemma for $P$.
Property (1) follows from the long exact cohomology sequence associated
to a short exact sequence of sheaves and the fact that tensoring with
an invertible sheaf is an exact functor. Property (2) follows since
$H^p(X, -)$ is an additive functor.

\medskip\noindent
To see (3) let $i : Z \to X$ be a reduced closed subspace with $|Z|$
irreducible. Let $i' : Z' \to Y$ and $f' : Z' \to Z$ be as in
Lemma \ref{lemma-finite-morphism-Noetherian} and set
$\mathcal{G} = f'_*\mathcal{O}_{Z'}$. We claim that
$\mathcal{G}$ satisfies properties (3)(a) and (3)(b) of
Lemma \ref{lemma-property-higher-rank-cohomological}
which will finish the proof. Property (3)(a) we have seen in
Lemma \ref{lemma-finite-morphism-Noetherian}. To see (3)(b) let
$\mathcal{I}$ be a nonzero quasi-coherent sheaf of ideals on $Z$.
Denote $\mathcal{I}' \subset \mathcal{O}_{Z'}$ the quasi-coherent
ideal $(f')^{-1}\mathcal{I} \mathcal{O}_{Z'}$, i.e., the
image of $(f')^*\mathcal{I} \to \mathcal{O}_{Z'}$.
By Lemma \ref{lemma-affine-morphism-projection-ideal} we have
$f_*\mathcal{I}' = \mathcal{I} \mathcal{G}$. We claim the common value
$\mathcal{G}' = \mathcal{I} \mathcal{G} = f'_*\mathcal{I}'$
satisfies the condition expressed in (3)(b).
First, it is clear that the support of $\mathcal{G}/\mathcal{G}'$
is contained in the support of $\mathcal{O}_Z/\mathcal{I}$
which is a proper subspace of $|Z|$ as $\mathcal{I}$ is a
nonzero ideal sheaf on the reduced and irreducible algebraic space $Z$.
Recall that $f'_*$, $i_*$, and $i'_*$ transform coherent modules into
coherent modules, see Lemmas \ref{lemma-finite-pushforward-coherent} and
\ref{lemma-i-star-equivalence}.
As $Y$ is a scheme and $\mathcal{L}$ is ample
we see from Lemma \ref{lemma-vanshing-gives-ample}
that there exists an $n_0$ such that
$$
H^p(Y, i'_*\mathcal{I}' \otimes_{\mathcal{O}_Y} f^*\mathcal{L}^{\otimes n}) = 0
$$
for $n \geq n_0$ and $p > 0$. Now we get
\begin{align*}
H^p(X, i_*\mathcal{G}' \otimes_{\mathcal{O}_X} \mathcal{L}^{\otimes n})
& =
H^p(Z, \mathcal{G'} \otimes_{\mathcal{O}_Z} i^*\mathcal{L}^{\otimes n}) \\
& =
H^p(Z,
f'_*\mathcal{I}' \otimes_{\mathcal{O}_Z} i^*\mathcal{L}^{\otimes n})) \\
& =
H^p(Z, f'_*(\mathcal{I}' \otimes_{\mathcal{O}_{Z'}}
(f')^*i^*\mathcal{L}^{\otimes n})) \\
& =
H^p(Z, f'_*(\mathcal{I}' \otimes_{\mathcal{O}_{Z'}}
(i')^*f^*\mathcal{L}^{\otimes n})) \\
& =
H^p(Z',
\mathcal{I}' \otimes_{\mathcal{O}_{Z'}} (i')^*f^*\mathcal{L}^{\otimes n})) \\
& =
H^p(Y, i'_*\mathcal{I}' \otimes_{\mathcal{O}_Y} f^*\mathcal{L}^{\otimes n}) = 0
\end{align*}
Here we have used the projection formula and the Leray spectral sequence
(see Cohomology on Sites, Sections
\ref{sites-cohomology-section-projection-formula} and
\ref{sites-cohomology-section-leray})
and Lemma \ref{lemma-finite-higher-direct-image-zero}.
This verifies property (3)(b) of
Lemma \ref{lemma-property-higher-rank-cohomological} as desired.
\end{proof}

\end{document}
